\section{Conclusion}
\label{sec:conclusion}

For object-level authorization, \mbox{MST-wi}'s predicates do not know whether
the acting subject is entitled to the object; they know what the interface, as
crawled, permitted. Instantiating them as defined produced three findings
(Section~\ref{sec:findings}).

\paragraph{Finding 1 (temporal)} As designed, the adjudication false-proofs on
legitimately granted access, deterministically (\ResObsOobVtwoFpRate{}
constructed; \ResRsMrBookkeepingFpRate{} Gitea), because an out-of-band grant fits
only the wrong one of the predicate's two representations. On the real system the
mismatch appears only under the crawl-time reading
(\ResRsMismatchSnapshot{}/\ResRsMismatchDenom{}, vs.\ \ResRsMismatchLive{} live),
with an endogenous control holding final visibility fixed: within this controlled
design the only manipulated explanatory variable is when the model was read, not
what was visible.

\paragraph{Finding 2 (oracle)} The natural repair---a readable, sound
authorization oracle that fully resolves the proposition under test---works, and
in that regime the evaluated metamorphic step
provides no observed incremental adjudication benefit over a direct consumer of the same oracle and outcome:
an empirical premise, formalised as a conditional (Proposition~\ref{lem:oracle-suff}), whose empirical face is that
a single ordinary request yields the same false-proof count (false proofs
\ResAbDirectObserveFP{} vs.~\ResAbVThreeFP{}), and that of the \ResAbNBola{}
true-violation records the paired comparison disagrees on only
\ResAbDisagreementN{}, each time in favor of the plain request; no margin was
pre-specified, so no inferential claim attaches.
On Gitea the comparison is limited to false-proof behavior, because no
vulnerability was present on the tested surface.

\paragraph{Finding 3 (measurement)} The reported specificity of
\ResLitReportedSpecPct{}\% is diagnosed as a specific published figure: it is
$\mathrm{TN}/(\mathrm{TN}+\mathrm{FP})$ over the follow-up inputs the crawl
could enumerate, with manually verified labels, and that denominator cannot
register a failure confined to a class the run did not exercise.

Support evidence: reading the implementation confirmed the semantics and added
the mechanism---the crawl model is frozen by a cache never invalidated, and a
missing crawl defaults to ``unreachable''---and the positive control is a sanity check that closes the loop: on a seeded
vulnerable target the surface-gated adjudicator detects all planted instances of
the fault pattern it was built for with zero false positives, and on four seeded
defect classes it detects the read bypass and nothing else---the write bypass
breaks a channel the relation does not test, and the inheritance and escalation
faults falsify the antecedent because the subject can already see the object.

\paragraph{What the two added experiments carry} The temporal finding is not
confined to one product. Under the same pre-registered protocol, frozen before
the second run was observed, GitLab CE \GlVersion{} commits \TmpGPolicyEvents{}
transitions over \TmpGCells{} cells, with entitlement and effect agreeing on
\TmpGEeqY{} of them: \TmpGStaleFP{} stale false proofs, all \TmpGFixed{} removed
by a recrawl with the same credentials, \TmpGLiveFP{} live false proofs and
\TmpGStillMiss{} residual misses. Pooled with Gitea this is \TmpBothCells{}
cells, \TmpBothEeqY{} agreements, \TmpBothStaleFP{} stale false proofs and
\TmpBothFixed{} fixes. The replication is of a protocol, not of a topology:
GitLab's project-membership surface answers 200 or 404, so it errs only by
omission---in \TmpGOracleDisagree{} cells it disagrees with the ledger, every one
of them a false deny, concentrated in the public
(\TmpGOracleDisagreePublic{}) and site-administration
(\TmpGOracleDisagreeAdmin{}) channels---and both properties are read from the
same run rather than one being derived from the other. The entitlement/effect
gate is in turn checked against a real defect: on a publicly fixed
object-versus-path ownership flaw, Gitea \CveVulnVersion{} yields
\CveTruePositives{} true positive with \CveFalseProofs{} false proofs, and the
fixed \CveFixedVersion{} yields \CveFixedAlarms{} alarms over \CveControls{}
control probes. That is a check of the gate on one defect we located and
supplied; it is not a detection rate.

\paragraph{What we do not claim, and how far the claims reach} We do not claim
that \mbox{MST-wi} is wrong or say anything about its accuracy as a detector; the
native run establishes a qualitative outcome, for one gate family, on one real
system that on the tested endpoints and transitions enforces the expected
decisions. Three ceilings stack and bound everything above. First, our two
relations are two of a catalogue of \ResLitMrCatalog{}, and both belong to one
reachability-gated family. Second, even the execution of the published relations
on the native engine reaches only that family: \NtMrFired{} of the \NtMrTotal{}
authorization relations fire (\NtMrExecuted{} executed; \NtKeyMrNames{}), and the content predicate
\texttt{userCanRetrieveContent} is never invoked (\NtUcrcCalls{} calls), so that
evidence is bounded the same way. Third, the real-system study exercises the
\emph{false-proof} face on a single constructed fault pattern, not the detection
face against a population of real flaws, and the Gitea channel proportions are
topology design parameters, not field prevalence. Together these fix the ceiling
on what a practitioner may transfer.

\paragraph{What a practitioner should do differently} Three operational
consequences follow, each bounded by what was measured. Do not read a
reachability-gated relation's alarm as proof: because the reading time is part of
the specification, an alarm raised against a crawl-time snapshot---rather than
against live state---cannot be promoted to \texttt{violation-proven} on the
relation's word alone. If a readable, sound authorization oracle is in fact
available to the analyst, consult it directly and drop the relation: for the
relations studied it adds no adjudication information once that antecedent holds
(Proposition~\ref{lem:oracle-suff}, Section~\ref{sec:stratum}). And treat a
reported input-level specificity as a corpus statistic rather than a guarantee of
scenario-level resolution: it can sit near one while a whole class of legitimate
access is misreported (Finding~3). These are not a claim that metamorphic testing
is worthless; they locate its contribution where no oracle exists, and bound it
by what the reachability reading can see.

\paragraph{What a stronger claim would require} Moving from characterization to something
stronger needs, first, a population of real systems that \emph{contain} the
flaw, so detection can be measured---neither the four seeded classes of
Section~\ref{sec:multiclass} nor the single publicly fixed defect of
Section~\ref{sec:cvepair} supplies that population: the first are planted in an
emulator, the second is a real defect replayed under a controlled topology, and
neither is a sample of defects found in the field; and second, relations whose gate is not the absence of
evidence---checkable against a subject's entitlements without pretending that
entitlements are reachable facts. The transferable condition---we call it the
\emph{same-observable condition}---is the one that killed our own
mechanism: a component's necessity must be tested against a substitute that
consumes the \emph{same} observable quantity as the ground truth; a substitute
that consumes less does not establish incremental value. The \mbox{MST-wi} case
study is an instance of that condition, not its scope: a component's necessity is
not established unless a fair control that lacks the component but consumes the
same decision-relevant observables is empirically or formally shown to be
weaker.
